\documentclass[11pt]{article}
\usepackage{enumitem}
\usepackage{latexsym}
\usepackage{amsfonts}
\usepackage{amsmath,amssymb,amsthm}
\usepackage{xcolor}

\usepackage{tikz}
\usetikzlibrary{arrows}
\usetikzlibrary{automata, positioning}

\setlength{\textheight}{8.5in}
\setlength{\textwidth}{6.0in}
\setlength{\headheight}{0in}
\addtolength{\topmargin}{-.5in}
\addtolength{\oddsidemargin}{-.5in}

\input{preamble.tex}

\newcommand{\solution}[1]{\paragraph{Solution}  }
\newcommand{\bl}[1]{\textcolor{blue}{#1}}
\newcommand{\rd}[1]{\textcolor{red}{#1}}

\begin{document}
\hw{1: Regular Languages}{1/08/2025}{}{Due:1/15/2025}

You should submit a typeset or \emph{neatly} written pdf on Gradescope.  The grading TA should not have to struggle to read what you've written; if your handwriting is hard to decipher, you will be asked to typeset your future assignments. Three bonus points if you use \LaTeX, and our template. You may collaborate with other students, but any written work should be your own. 

\begin{enumerate}

    \item One perspective of this class is that we are developing a theory of representation, or definability. For finite sets, this is easy, but it becomes more challenging for infinite languages. This class is really about the ability of different finite objects to understand and describe infinite sets.

    \begin{enumerate}
        \item Prove every finite language is regular.
        
            \begin{answer}

            \end{answer}
            
        \item Suppose we relax our definition of DFA too much. We define a DIA\footnote{Deterministic Infinite Automata} to be exactly like a DFA except $|Q|=|\mathbb{N}|$, and $\delta, F$ are adjusted appropriately. Prove that \textit{every} language $L \subseteq \Sigma^*$ is decided/accepted by a DIA.
        
            \begin{answer}

            \end{answer}
            
    \end{enumerate}
    
    %https://drawautomata.xyz/ may help you in constructing diagrams if you choose to submit latex. 
    \item Give a DFA for the following languages. Your solution doesn't have to be minimal, but it may assist the TA in grading. Here $int(w,2)$ converts the binary string $w$ to a number. Assume that $\Sigma= \{0,1\}$. 
    \begin{enumerate}

        \item $L_1 = \{1^n ~|~ n \equiv 3 \pmod 4\}$

            \begin{answer}

            \end{answer}
            
        \item $L_2 = \{w \in \Sigma^* ~|~ int(w,2) \equiv 3 \pmod 4\}$

            \begin{answer}
        
            \end{answer}
            
        \item $L_3 = \{w \in \Sigma^* ~|~ int(w,2) \equiv 3 \pmod 5\}$

            \begin{answer}
        
            \end{answer}
            
    \end{enumerate}
    
    \item Let $L$ be a language from the alphabet $\Sigma = \{a,b\}$. Let $L'$ be a language over the alphabet $\Sigma' = \{a,b,c\}$ where $$L' = \{xcy \in \Sigma'^*~|~ xy \in L\}$$ We take the strings in $L$, and add a single new symbol $c$ anywhere in the string, even possibly at the beginning or end. Prove that if $L$ is regular, then so is $L'$.
        
        \begin{answer}
            
        \end{answer}
        
    \item Define $L(a,b) = \{1^{an + b}~|~ n \in \mathbb{N}\}$. Let $L \subseteq 1^*$. Prove that $L$ is regular if and only if there exists numbers $a_1,b_1,...,a_k,b_k \in \mathbb{N}$ and a finite language $F\subseteq 1^*$ such that $$L = L(a_1,b_1) \cup ... \cup L(a_k,b_k) \cup F$$ You are proving the unary regular languages are exactly those which are unions of arithmetic progressions, and some finite extra strings.
        
        \begin{answer}
        
        \end{answer}
        
\end{enumerate}
\end{document}